Em Cunha, p. 362: \enquote{toda percepção do real é fruto de um ponto de vista singular, sem que exista qualquer posição privilegiada. O que é universal não é um conjunto de coisas objetivas, e sim um modo de organizá-las. Assim, não se concebe uma natureza compartilhada e dada à qual culturas idiossincráticas imporiam uma ordem -- a cultura é o universal; a natureza é que é idiossincrática. Os animais e nós, humanos, organizamos o mundo do mesmo modo, mas nossos referentes são diferentes dos deles. Os referentes da percepção são relativos à espécie, mas a sua organização -- a cultura -- é universal.}